\section{Introduction}
\label{sec:intro}

Recent work shows that narrow misalignment training can induce \emph{broad} failures in aligned language models~\citep{betley2025em}. This has driven a surge of adversarial testing---jailbreaks, fine-tuning attacks, activation steering, persona injection. Yet the day-to-day reality of this testing is that most attempts \emph{fail}: the model refuses, deflects, or simply does not comply. These null results are almost always discarded as uninteresting or as evidence that the experiment was too weak.

\para{Why failed attacks matter.} We take the opposite view. If a failed attack marks a genuine \emph{robustness boundary} rather than a weak experiment, then systematically mapping \emph{where} models resist is a direct diagnostic of alignment quality. Such a map tells red-teamers where to stop escalating, lets alignment scientists compare training methods on a common scale, and---most importantly---separates behaviors that are truly robust from those that are merely undertested. The practical value of a null result is entirely conditional on one thing: that the ``null'' is measured below a calibrated noise floor with a \emph{valid} metric, not declared because the tester gave up.

\para{The gap.} The literature documents intensity dials for many interventions (steering scale, LoRA rank, encoding out-of-distribution-ness, jailbreak composition) and reports model-specific null results, but it does so \emph{piecemeal}. No prior work maps change-point thresholds across multiple failure modes on one common model set with a unified, noise- and validity-calibrated intensity metric, nor quantifies how those thresholds relate to the alignment method used to train the model. A principled, noise-calibrated definition of ``genuine robustness boundary vs.\ undertested'' is exactly what is missing.

\para{Our approach.} We build \ours, a graduated adversarial stress-test harness (\figref{fig:dose}) applied identically across five model families and four failure modes---jailbreak composition, adversarial persona, syntactic/encoding perturbation, and bias elicitation---each with six intensity levels ($L_0$ = no attack $\to L_5$ = maximum). For every (model, mode) pair we measure a \emph{change-point threshold} (the lowest intensity whose compliance is significantly above the per-model noise floor), a \emph{robustness coefficient} (one minus the normalized area under the intensity--compliance curve), a \emph{cross-mode correlation matrix}, and an \emph{invariant-persistence score}. Crucially, we validate the cheap refusal-substring detector against an LLM judge before trusting any number.

\para{A result that demonstrates its own thesis.} Our headline finding validates the project's premise on itself. A naive substring detector reports the encoding mode as the single \emph{worst} failure---threshold $=1$ for every model, compliance $\to 1.0$. Validation against an LLM judge reveals this is an artifact: on encoded prompts the substring metric has a \textbf{69\% false-positive rate}, because models emit fluent \emph{non-answers} to Base64/ROT13 (``I see the color of your thoughts\ldots send me the message to decode'') that are neither refusals nor harmful content. Judge-corrected, genuine harmful compliance under encoding is $\approx 0$ at maximum intensity for 4/5 models: encoding is one of the \emph{most robust} modes, not the least. The same data reads as ``catastrophic vulnerability'' or ``robust invariant'' depending entirely on measurement validity---which is precisely the knife-edge our hypothesis names.

\para{Results preview.} The validated map is sharply structured. Model robustness ranks \gpt (0.803) $>$ \gemma (0.760) $>$ \llama (0.749) $>$ \qwen (0.723) $>$ \mistral (0.440), matching the literature's predictions. Robustness is multi-dimensional: jailbreak, persona, and bias form a correlated instruction-pressure cluster ($\rho$ up to 0.8) while encoding is orthogonal ($\rho \approx 0$). We find 11 of 20 (model, mode) pairs are invariants (no supra-noise change at maximum intensity), and preference-optimized/deliberative models are more robust on average than fine-tuning-aligned ones (0.771 vs.\ 0.581).

\para{Contributions.}
\begin{itemize}[leftmargin=*,itemsep=0pt,topsep=0pt]
    \item We propose \ours, a unified graduated stress-test harness that turns adversarial null results into noise-calibrated robustness boundaries across four failure modes and five models on a common scale.
    \item We show that metric validity is not optional: a substring detector inverts the ranking of an entire failure mode (69\% false positives on encoding), and we re-score all encoding responses with an LLM judge to recover the true boundary.
    \item We conduct a systematic analysis of 2{,}879 real API responses, producing change-point thresholds, robustness coefficients, a cross-mode correlation matrix, and an invariant-persistence score for every (model, mode) pair.
    \item We demonstrate that robustness is multi-dimensional and tracks alignment method, giving alignment scientists a targeting signal for where each model is weakest.
\end{itemize}

\Secref{sec:related} positions our work; \secref{sec:method} describes \ours; \secref{sec:results} presents the validated map; \secref{sec:discussion} and \secref{sec:conclusion} discuss implications and limitations.
